\documentclass[11pt]{article}
\usepackage[margin=1in]{geometry}
\usepackage{amsmath,amssymb}
\usepackage{booktabs}
\usepackage{hyperref}

\title{EEG-fMRI in awake rat and whole-brain simulations show decreased brain responsiveness to sensory stimulations during absence seizures}
\author{}
\date{}

\begin{document}
\maketitle

\begin{abstract}
In patients with absence epilepsy, recurring seizures decrease quality of life and lead to untreatable comorbidities. Absence seizures are characterized by spike-and-wave discharges (SWDs) on EEG and a transient alteration of consciousness; how the brain responds to external stimuli during and outside seizures remains unclear. We investigated responsiveness to visual and somatosensory stimulation in GAERS, a rat model of absence epilepsy. Animals were imaged awake (non-curarized) using a quiet zero-echo-time (ZTE) fMRI sequence. Whole-brain hemodynamic responses to visual and whisker stimulation were compared between interictal and ictal states. A mean-field simulation model explained neural responsiveness changes. Whole-brain responses to both sensory stimulations were suppressed during seizures; cortical hemodynamic responses were negatively polarized despite applied stimuli. The mean-field simulation revealed restricted propagation of activity, agreeing with fMRI findings.
\end{abstract}

\section{Introduction}
Absence seizures appear in children aged 5--7 years, involve a brief impairment of consciousness and unresponsiveness, and feature a regular 2--5\,Hz SWD pattern on EEG \citep{loiseau1995,szaflarski2010}. Untreated, seizures can occur from a few to hundreds of times per day, with impairments of attention, memory, and mood \citep{crunelli2020}. Human neuroimaging suggests lack of conscious information processing due to impaired fronto-parietal networks, arousal systems in thalamus and brainstem, or the default mode network \citep{blumenfeld2012,luo2011}. Most existing imaging studies are at rest, without external stimulation.

GAERS (Genetic Absence Epilepsy Rats from Strasbourg) is a well-established rat model of absence epilepsy \citep{depaulis2016}. Absence seizures in GAERS initiate in the barrel field of the somatosensory cortex \citep{david2008,meeren2002,polack2007}, and propagate to thalamus \citep{huguenard2019}. During SWDs, majority of frontoparietal cortical and thalamocortical relay neurons decrease firing \citep{mccafferty2023}. fMRI shows BOLD decreases in parietal and occipital cortex but increases in subcortical regions \citep{david2008,mccafferty2023}. We used EEG-fMRI in awake GAERS, with visual and whisker stimulation, and a whole-brain Virtual Brain mean-field simulation \citep{divolo2019} to describe interictal-vs-ictal responsiveness.

\section{Materials and Methods}
\subsection{Animals}
Eleven adult 8--12 month-old GAERS rats were used (260 $\pm$ 21\,g; 6 males, 5 females). All experiments were approved by the local animal welfare committee (Comit\'e Local GIN, C2EA-04) per EU Directive 2010/63/EU.

\subsection{EEG implantation and restraint}
Carbon fibre electrodes were implanted over the right motor cortex (AP +2, ML +2.5\,mm) and right primary somatosensory cortex (AP $-2.5$, ML +3\,mm); the reference/ground electrode was placed on the cerebellum (AP $-12$, ML +2\,mm). Rats were habituated to restraint and ZTE gradient noise for 15 to 60 minutes per day for 8 days.

\subsection{EEG-fMRI acquisition}
MRI was acquired at 9.4\,T (Bruker Biospec) with an in-house transmit-receive loop coil. Anatomical: T1-FLASH (TR 530\,ms, TE 4\,ms, flip 18$^\circ$, matrix $128\times128$, 51 slices, FOV $32\times32$\,mm$^2$, resolution $0.25\times0.25\times0.5$\,mm$^3$). Functional: 3D ZTE (TR 0.971\,ms, TE 0\,ms, flip 4$^\circ$, pulse 1\,$\mu$s, bandwidth 150\,kHz, oversampling 4, matrix $60\times60\times60$, FOV $30\times30\times60$\,mm$^3$, resolution $0.5\times0.5\times1$\,mm$^3$, polar undersampling 5.64, 2060 projections, $\sim$2\,s per volume, 1350 volumes per 45-min scan).

Visual stimulation: blue LED (470\,nm) at 3\,Hz, 6\,s blocks, pulse length 166\,ms. Whisker stimulation: air-puff 2\,Hz, 6\,s blocks, pulse length 250\,ms, pressure 2.3\,bar. Stimulations were manually triggered during ictal or interictal states using live EEG. EEG: 1024\,Hz, 50\,Hz notch and 1--90\,Hz Butterworth band-pass filters, with TTL triggers logged per stimulation onset.

\subsection{Sessions}
Visual group: 14 sessions across 5 rats. Whisker group: 14 sessions across 4 rats. Three of 11 animals (27\%) were excluded due to noisy EEG; one (1\%) due to lost implant. 6/28 sessions (21\%) were excluded for motion, 4 (14\%) for no seizures, and 1 (4\%) for air-puff failure.

\subsection{Simulation}
Each brain region was modelled with an Adaptive Exponential LIF (AdEx) mean-field \citep{divolo2019}. The model parameterises asynchronous irregular (AI, interictal) or SWD (ictal) dynamics by the strength of an adaptation current. A connectivity matrix of 496 rat-brain regions (BAMS rat connectome \citep{bota2012}) in The Virtual Brain platform \citep{sanzleon2015} simulated stimulus propagation.

\subsection{Statistics}
Extreme beta-values in anatomical ROIs from SPM were compared with a two-sample $t$-test. F-contrast maps were corrected for multiple comparisons by cluster-level correction ($p<0.05$). Statistical power calculation: with $\alpha=0.05$, power 0.8, matched pairs, we detect effect size 0.37 with 4 animals (4 sessions per animal $\times$ 11 seizure/control pairs per session).

\section{Results}
\subsection{fMRI setup in awake rats}
Spatial SNR was $\sim$25; temporal SNR 30--60 in the brain. Gradient-echo-EPI produced 1--1.5\% signal change vs.\ $\sim$0.3\% for ZTE, but ZTE had less susceptibility distortion and less EEG artefact. Peak acoustic noise: ZTE 78.7\,dB vs.\ EPI 114.5\,dB (ZTE 35.8\,dB quieter, $\sim$62$\times$ weaker sound pressure). Mean framewise head translation 0.69\,$\pm$\,0.32\,$\mu$m; head rotation 0.79\,$\pm$\,0.55$^\circ$. Maximum displacement 12.8\,$\pm$\,11.6\,$\mu$m (0.03\,$\pm$\,0.02 voxels); maximum rotation 18.1\,$\pm$\,12.9$^\circ$. Motion occurrences: 0.43\,$\pm$\,0.45 motions/min (current) vs.\ 1.0\,$\pm$\,0.20 motions/min (previous EPI), 0.48\,$\pm$\,0.23 motions/min (MB-SWIFT). Respiration: 2.2\,$\pm$\,0.4\,Hz (visual) and 2.0\,$\pm$\,0.3\,Hz (whisker).

\subsection{Response to sensory stimuli}
During interictal visual stimulation, activations ($p<0.05$, cluster-corrected) were seen in visual cortex, superior colliculus, thalamus (including lateral geniculate nucleus), and frontal cortex (prelimbic, cingulate, secondary motor). During ictal visual stimulation, visual cortex responses were less pronounced while superior colliculus responses remained stable. Interictal-vs-ictal comparison: 136\% more active voxels during interictal vs.\ ictal in the visual group. During interictal whisker stimulation, responses were seen in somatosensory cortex including barrel field, auditory cortex, thalamus including ventral posteromedial nucleus, and medial frontal cortex. Whisker group had 179\% more active voxels during interictal vs.\ ictal.

Extreme beta-values in the visual cortex ROI: baseline $2.8 \pm 1.7$; seizure $-6.0 \pm 2.0$ ($p<0.001$). Stimulation ending a seizure: $8.1 \pm 7.0$ vs.\ $-6.0 \pm 2.0$ ($p<0.001$). Barrel cortex ROI: baseline $4.1 \pm 1.9$; seizure $-9.0 \pm 1.9$ ($p<0.001$); stimulation ending a seizure: $4.8 \pm 2.9$ vs.\ $-9.0 \pm 1.9$ ($p<0.001$). Pooled seizure effect (visual + whisker): significant cortical and subcortical changes over 895 voxels; negative cortical and positive deep-brain polarity.

\begin{table}[h]
\centering
\caption{Extreme beta-values by ROI and condition (mean $\pm$ SD).}
\begin{tabular}{lccc}
\toprule
ROI & Baseline & Seizure & Stimulation ended seizure \\
\midrule
Visual cortex & $2.8 \pm 1.7$ & $-6.0 \pm 2.0$ & $8.1 \pm 7.0$ \\
Barrel cortex & $4.1 \pm 1.9$ & $-9.0 \pm 1.9$ & $4.8 \pm 2.9$ \\
\bottomrule
\end{tabular}
\end{table}

\subsection{Mean-field simulation}
The AdEx mean-field model produced AI and SWD dynamics matching interictal and ictal states. SWD spikes correlated with periods of hyper-polarization driven by a strong adaptation current. Whole-brain simulations of a primary visual cortex stimulus showed drastically restrained propagation during ictal periods, while interictal periods showed large variation in regions linked to the stimulated area, consistent with fMRI observations.

\section{Discussion}
Zero echo-time fMRI of non-curarized awake rats enabled detection of absence seizures, with ZTE preferred over EPI for quieter acquisition and fewer susceptibility artefacts. ZTE fMRI is sensitive to changes in blood flow \citep{luh2000,restom2007} and has been suggested to be 67\% more sensitive than standard BOLD EPI.

Interictal responses reached visual cortex, superior colliculus, thalamus, and frontal cortex \citep{sefton2015,williams1999,muir1996}. Ictal cortical responses were negatively polarized, consistent with the ``wave'' (silence) component of SWD being dominant and with reduced neuronal activity \citep{fisher1977,gloor1978,inoue1993,kostopoulos1982,mccafferty2023}. When stimulation ended a seizure, fMRI response amplitudes were higher, suggesting a threshold effect that shifted the brain back to a responsive state. Simulation results agreed with fMRI, showing restricted stimulus propagation during ictal periods. These findings may explain cognitive comorbidities in absence epilepsy and inform therapeutic approaches.

\bibliographystyle{plain}
\bibliography{references}

\end{document}
